%% DAFTAR SINGKATAN -- singkatan yang digunakan dalam skripsi ini.
%%
%% Disusun dengan mengekstraksi setiap token huruf kapital dari Bab I-IV lalu
%% mempertahankan yang benar-benar merupakan singkatan. Sengaja dikecualikan: kata kunci
%% Cypher dan label graf (MERGE, SUBJECT, SAME_CONCEPT), awalan pengenal rekomendasi
%% (APT, BB, BPM, LIPID, MTRA...), pengenal konsep UMLS (C0151744), dan nama metrik (F1).
%%
%% Perluasan singkatan sengaja dipertahankan dalam bahasa Inggris karena itulah bentuk
%% resmi istilahnya; penjelasan tambahan diberikan dalam bahasa Indonesia.
%%
%% \crosscheck{Jalankan ulang ekstraksi setelah revisi besar; dan jika awalan pengenal
%% rekomendasi banyak dipakai pada tabel Bab IV, pertimbangkan menambahkan kunci singkat
%% tersendiri di sana alih-alih di sini.}

\begin{longtable}{llp{3.6in}}
	ACC     & \hspace{1cm} = & American College of Cardiology \\
	ACS     & \hspace{1cm} = & \emph{Acute Coronary Syndrome} \\
	AHA     & \hspace{1cm} = & American Heart Association \\
	API     & \hspace{1cm} = & \emph{Application Programming Interface} \\
	ATC     & \hspace{1cm} = & \emph{Anatomical Therapeutic Chemical} (klasifikasi obat WHO) \\
	BMI     & \hspace{1cm} = & \emph{Body Mass Index} \\
	BP      & \hspace{1cm} = & \emph{Blood Pressure} \\
	CABG    & \hspace{1cm} = & \emph{Coronary Artery Bypass Grafting} \\
	CALBC   & \hspace{1cm} = & \emph{Collaborative Annotation of a Large Biomedical Corpus} \\
	CCD     & \hspace{1cm} = & \emph{Chronic Coronary Disease} \\
	CDSS    & \hspace{1cm} = & \emph{Clinical Decision Support System} \\
	CIG     & \hspace{1cm} = & \emph{Computer-Interpretable Guideline} \\
	CKD     & \hspace{1cm} = & \emph{Chronic Kidney Disease} \\
	COR     & \hspace{1cm} = & \emph{Class of Recommendation} \\
	CSV     & \hspace{1cm} = & \emph{Comma-Separated Values} \\
	CUI     & \hspace{1cm} = & \emph{Concept Unique Identifier} (pengenal konsep UMLS) \\
	DAPT    & \hspace{1cm} = & \emph{Dual Antiplatelet Therapy} \\
	CRS     & \hspace{1cm} = & \emph{Clinical Recommendation System} \\
	EHR     & \hspace{1cm} = & \emph{Electronic Health Record} \\
	FHIR    & \hspace{1cm} = & \emph{Fast Healthcare Interoperability Resources} \\
	GDMT    & \hspace{1cm} = & \emph{Guideline-Directed Medical Therapy} \\
	GLIF    & \hspace{1cm} = & \emph{GuideLine Interchange Format} \\
	GPT     & \hspace{1cm} = & \emph{Generative Pre-trained Transformer} \\
	HDL     & \hspace{1cm} = & \emph{High-Density Lipoprotein} \\
	HL7     & \hspace{1cm} = & \emph{Health Level Seven} (organisasi standar interoperabilitas) \\
	ICD     & \hspace{1cm} = & \emph{International Classification of Diseases} \\
	IQR     & \hspace{1cm} = & \emph{Interquartile Range} \\
	IRB     & \hspace{1cm} = & \emph{Institutional Review Board} \\
	KG      & \hspace{1cm} = & \emph{Knowledge Graph} \\
	LDL-C   & \hspace{1cm} = & \emph{Low-Density Lipoprotein Cholesterol} \\
	LLM     & \hspace{1cm} = & \emph{Large Language Model} \\
	LOE     & \hspace{1cm} = & \emph{Level of Evidence} \\
	LOINC   & \hspace{1cm} = & \emph{Logical Observation Identifiers Names and Codes} \\
	LVEF    & \hspace{1cm} = & \emph{Left Ventricular Ejection Fraction} \\
	MACE    & \hspace{1cm} = & \emph{Major Adverse Cardiovascular Events} \\
	MeSH    & \hspace{1cm} = & \emph{Medical Subject Headings} \\
	MI      & \hspace{1cm} = & \emph{Myocardial Infarction} \\
	MIMIC   & \hspace{1cm} = & \emph{Medical Information Mart for Intensive Care} \\
	MRR     & \hspace{1cm} = & \emph{Mean Reciprocal Rank} \\
	NER     & \hspace{1cm} = & \emph{Named-Entity Recognition} \\
	NICE    & \hspace{1cm} = & National Institute for Health and Care Excellence (Britania Raya) \\
	NLP     & \hspace{1cm} = & \emph{Natural Language Processing} \\
	NSAID   & \hspace{1cm} = & \emph{Non-Steroidal Anti-Inflammatory Drug} \\
	NSTEMI  & \hspace{1cm} = & \emph{Non-ST-Elevation Myocardial Infarction} \\
	OAC     & \hspace{1cm} = & \emph{Oral Anticoagulant} \\
	OHDSI   & \hspace{1cm} = & \emph{Observational Health Data Sciences and Informatics} \\
	OMOP    & \hspace{1cm} = & \emph{Observational Medical Outcomes Partnership} \\
	OWL     & \hspace{1cm} = & \emph{Web Ontology Language} \\
	PCI     & \hspace{1cm} = & \emph{Percutaneous Coronary Intervention} \\
	PDE5i   & \hspace{1cm} = & \emph{Phosphodiesterase Type 5 Inhibitor} \\
	PHI     & \hspace{1cm} = & \emph{Protected Health Information} \\
	PMI     & \hspace{1cm} = & \emph{Pointwise Mutual Information} \\
	PMID    & \hspace{1cm} = & \emph{PubMed Identifier} \\
	RAG     & \hspace{1cm} = & \emph{Retrieval-Augmented Generation} \\
	RCT     & \hspace{1cm} = & \emph{Randomised Controlled Trial} \\
	RDF     & \hspace{1cm} = & \emph{Resource Description Framework} \\
	RO      & \hspace{1cm} = & \emph{Research Objective} (tujuan penelitian) \\
	RQ      & \hspace{1cm} = & \emph{Research Question} (pertanyaan penelitian) \\
	RxNorm  & \hspace{1cm} = & terminologi obat ternormalisasi dari \emph{US National Library of Medicine} \\
	SIHD    & \hspace{1cm} = & \emph{Stable Ischemic Heart Disease} (subhimpunan dari CCD) \\
	SNOMED~CT & \hspace{1cm} = & \emph{Systematized Nomenclature of Medicine -- Clinical Terms} \\
	STEMI   & \hspace{1cm} = & \emph{ST-Elevation Myocardial Infarction} \\
	STU3    & \hspace{1cm} = & \emph{Standard for Trial Use, Release 3} (versi FHIR sebelum R4) \\
	TG      & \hspace{1cm} = & \emph{Triglyceride} \\
	UMLS    & \hspace{1cm} = & \emph{Unified Medical Language System} \\
	URI     & \hspace{1cm} = & \emph{Uniform Resource Identifier} \\
	WHO     & \hspace{1cm} = & World Health Organization 
\end{longtable}

\vspace{0.6cm}

\noindent Notasi statistik:

\vspace{0.2cm}

\begin{longtable}{llp{3.6in}}
	$\kappa$ & \hspace{1cm} = & koefisien kesepakatan Cohen (\emph{Cohen's kappa}) \\
	$F_1$    & \hspace{1cm} = & rerata harmonik \emph{precision} dan \emph{recall} \\
	$I^2$    & \hspace{1cm} = & statistik heterogenitas pada meta-analisis 
\end{longtable}
